% answers_short.tex —— 2008 年·普通高等学校招生考试·安徽卷·文科 · 简短答案速查
% 由 tools/make_short_answers.py 自 answers.tex 的 \gksolhead/\gkanswer 提取，
% 并与 2008/2008_数学真题_安徽文_答案要点.md 逐题比对后生成，勿手改（改源 MD 后重跑）。

\begin{gknotice}
\textbf{使用说明}：本节只给最终答案便于快速核对，完整推导见解析版。答案为编写者独立推导，\textbf{非官方标准答案}。
\end{gknotice}

\gkpart{选择题}
\begin{center}\begin{tabular}{ccccccccccccc}
\toprule
题号 & 1 & 2 & 3 & 4 & 5 & 6 & 7 & 8 & 9 & 10 & 11 & 12 \\
\midrule
答案 & A & B & B & B & A & C & A & D & A & D & C & C \\
\bottomrule
\end{tabular}\end{center}

\gkpart{填空题}
\begin{itemize}[leftmargin=2.6em]
  \item[\textbf{13.}] $[3,+\infty)$。
  \item[\textbf{14.}] $4$。
  \item[\textbf{15.}] $-1$。
  \item[\textbf{16.}] $\dfrac{4\pi}3$。
\end{itemize}

\gkpart{解答题}
\begin{itemize}[leftmargin=2.6em]
  \item[\textbf{17.}] （1）$\pi$；（2）$\left[-\dfrac{\sqrt3}2,1\right]$。
  \item[\textbf{18.}] （1）$\dfrac{27}{1000}$；（2）$\dfrac{11}{60}$。
  \item[\textbf{19.}] （1）$60^\circ$；（2）$\dfrac23$。
  \item[\textbf{20.}] （1）$a=1$；（2）$[-2,0]$。
  \item[\textbf{21.}] （1）$c=1$ 时 $a_n=a$；$c\neq1$ 时 $a_n=1+(a-1)c^{n-1}$；（2）$S_n=2-\dfrac{n+2}{2^n}$；（3）证明见解析。
  \item[\textbf{22.}] （1）$\dfrac{x^2}8+\dfrac{y^2}4=1$；（2）证明见解析；（3）$\dfrac{16\sqrt2}3$。
\end{itemize}
